\documentclass[11pt,a4paper]{article}
\usepackage[utf8]{inputenc}
\usepackage{geometry}
\geometry{a4paper, margin=1in}
\usepackage{hyperref}
\usepackage{graphicx}
\usepackage{listings}
\usepackage{xcolor}
\usepackage{amsmath}
\usepackage{booktabs}
\usepackage{float}
\usepackage{tikz}
\usepackage{eso-pic}
\usepackage{setspace} % Added for line spacing
\usepackage{parskip}  % Added for better paragraph spacing
\usepackage[T1]{fontenc}
\usepackage{microtype}
\usepackage[hyphens]{url}
\usepackage{caption}
\usepackage{subcaption}
\usepackage{titlesec}
\usepackage{sectsty}
\usepackage{fancyhdr}
\usepackage{enumitem}

% Allow more flexibility in line breaks to prevent border crossing
\emergencystretch 3em
\sloppy

% --- Nice Border for All Pages ---
\AddToShipoutPictureBG{%
    \begin{tikzpicture}[remember picture,overlay]
        % Outer thick border
        \draw[line width=1.5pt, rounded corners=10pt]
            ([shift={(1.2cm,-1.2cm)}]current page.north west)
            rectangle
            ([shift={(-1.2cm,1.2cm)}]current page.south east);
        % Inner thin border
        \draw[line width=0.5pt, rounded corners=8pt]
            ([shift={(1.35cm,-1.35cm)}]current page.north west)
            rectangle
            ([shift={(-1.35cm,1.35cm)}]current page.south east);
    \end{tikzpicture}%
}

% ------------------------- Colors & Formatting ------------------------- %
\definecolor{primary}{HTML}{1F3A60}
\definecolor{secondary}{HTML}{4B6A9B}
\definecolor{accent}{HTML}{E74C3C}
\definecolor{lightgray}{HTML}{F7F9FA}
\definecolor{border}{HTML}{D1D9E0}

\sectionfont{\color{primary}\Large\bfseries}
\subsectionfont{\color{secondary}\large\bfseries}

% Remove header/footer globally
\pagestyle{empty}
\renewcommand{\headrulewidth}{0pt}

\begin{document}

\begin{center}
    {\Huge \textbf{Machine Learning Assignment 1}} \\[0.5cm]
    {\LARGE \textbf{Polynomial Regression Report}} \\[0.5cm]
    {\Large \textbf{Rigorous Data-Driven Analysis of Phase I \& II}} \\[0.5cm]
    \textbf{Name: Sushmit Biswas \quad | \quad Roll Number: BT2024038}
\end{center}
\vspace{0.5cm}
\hrule height 1pt
\vspace{0.8cm}

\textbf{Abstract:} This report documents a data-driven approach to establishing optimized Polynomial Regression models for two personalized datasets. The methodology utilizes exploratory data analysis (EDA) to evaluate variable correlations, followed by 5-Fold Cross Validation to programmatically identify robust polynomial degrees without introducing external regularization. By evaluating Out-of-Fold (OOF) predictions, models achieving out-of-fold $R^2$ estimates above $0.90$ were successfully constructed.

% ------------------------- Introduction ------------------------- %
\section{Analytical Premise}
Standard polynomial approximations rely on $w_0 + \sum_{i=1}^{d} w_i \phi_i(\mathbf{x})$. Blindly setting $d$ without structural evaluation exposes the model to severe underfitting or Runge-like oscillation overfitting. To reduce the risk of overfitting, polynomial degrees within this project were selected based strictly on validation-fold mean squared error using 5-Fold Cross Validation (Shuffled, random seed=42). This ensures the reported metrics reflect generalization capabilities rather than overfitting. All available input features were retained, allowing the polynomial expansion to represent nonlinear and interaction effects.

\section{Phase 1: Steam Turbine Optimization (var1)}
The goal of Phase 1 is estimating the Net Power Score ($y$) based on physical turbine variables.

\subsection{Inclusive Feature Representation}
We compute an inclusive Pearson Correlation matrix on all given dataset features $\{x_1, \dots, x_6\}$. 
\begin{figure}[htbp]
    \centering
    \includegraphics[width=0.65\textwidth]{images/phase1_heatmap.png}
    \caption{Pearson Correlation Heatmap displaying relational bounds in Phase 1.}
\end{figure}
All available features were retained because weak marginal correlation does not rule out nonlinear or interaction effects in a polynomial model.

\subsection{Cross-Validation and Polynomial Sizing}
Using all six available features, five-fold cross-validation evaluated candidate degrees from 1 through 10.
\begin{figure}[htbp]
    \centering
    \includegraphics[width=0.8\textwidth]{images/phase1_cv_curve.png}
    \caption{Five-fold cross-validation MSE across candidate polynomial degrees.}
\end{figure}
At \textbf{Degree 3}, the validation error curve reaches an optimal minimum among degrees 1--10 before increasing degree complexity causes overfitting on validation sets. From degree 6 onward, the number of polynomial terms becomes comparable to or larger than the number of observations, causing near-zero training error while validation error remains substantially higher. Therefore, validation MSE rather than training MSE was used for degree selection.

\subsection{Phase 1 Empirical Results}
\begin{itemize}[noitemsep]
    \item \textbf{Features Used:} $x_1, x_2, x_3, x_4, x_5, x_6$
    \item \textbf{Selected Degree:} 3
    \item \textbf{Training MSE:} $0.7727$, \textbf{Training $R^2$:} $0.9272$
    \item \textbf{Out-Of-Fold Validation MSE:} $0.9895$
    \item \textbf{Out-Of-Fold Validation $R^2$:} $0.9067$
\end{itemize}
These values indicate that the selected polynomial captures a substantial portion of the observed target variation. Residuals are broadly centered around zero, with some outliers and no dominant visible structure.
\begin{figure}[htbp]
    \centering
    \includegraphics[width=0.6\textwidth]{images/phase1_residuals.png}
    \caption{Residual distribution mapping relative predictive error across constraints.}
\end{figure}


% ------------------------- Phase 2 ------------------------- %
\section{Phase 2: Subterranean Thermal Mapping (var2)}

\subsection{Automated Structural Mapping}
Spatial relationships inherently harbor latent collinearities. 
\begin{figure}[htbp]
    \centering
    \begin{subfigure}{0.48\textwidth}
        \centering
        \includegraphics[width=\textwidth]{images/phase2_pairplot.png}
        \caption{Sensor Pair-Plot Topology.}
    \end{subfigure}
    \hfill
    \begin{subfigure}{0.48\textwidth}
        \centering
        \includegraphics[width=\textwidth]{images/phase2_heatmap.png}
        \caption{Matrix Variance Correlates.}
    \end{subfigure}
    \caption{Feature evaluations confirming mathematical retention bounds.}
\end{figure}
Again, all three structural coordinates ($x_1, x_2, x_3$) were retained to preserve the multi-dimensional prediction topology against potential non-linear dependencies.

\subsection{Degree Optimization Boundaries}
A logarithmic sweep mapping error arrays from degree 1 up to polynomial explosion bounds was evaluated. Validation MSE decreases through degree 8 and then increases sharply, while training MSE continues to decrease. This gap indicates overfitting and numerical instability at higher degrees.
\begin{figure}[htbp]
    \centering
    \includegraphics[width=0.7\textwidth]{images/phase2_cv_curve.png}
    \caption{Five-fold cross-validation MSE across candidate polynomial degrees.}
\end{figure}
The minimal automated Validation MSE confirms an optimal fit at \textbf{Degree 8} among candidate degrees 1--20.

\subsection{Phase 2 Empirical Results}
\begin{itemize}[noitemsep]
    \item \textbf{Features Used:} $x_1, x_2, x_3$
    \item \textbf{Discovered Degree:} 8
    \item \textbf{Training MSE:} $0.1731$, \textbf{Training $R^2$:} $0.9965$
    \item \textbf{Out-Of-Fold Validation MSE:} $0.2629$
    \item \textbf{Out-Of-Fold Validation $R^2$:} $0.9947$
\end{itemize}
\begin{figure}[htbp]
    \centering
    \includegraphics[width=0.55\textwidth]{images/phase2_true_vs_pred.png}
    \caption{True output mappings accurately predicting geological boundaries.}
\end{figure}

\section{Computational Complexity \& Overparameterization}
During the exhaustive degree search, models were evaluated up to the maximum candidate degrees specified for this analysis (Degree 10 for Phase 1 and Degree 20 for Phase 2). It is mathematically critical to note that high-degree polynomial expansions trigger severe overparameterization. For instance, evaluating Phase 1 at Degree 10 across 6 features generates over 8,000 distinct polynomial interaction parameters. When fitted against approximately 1,000 training observations, the design matrix becomes pathologically rank-deficient (where $P > N$).

This disparity causes the standard Ordinary Least Squares (OLS) estimator to produce nearly zero training error. However, this high-dimensional polynomial surface exhibits wild mathematical oscillations between known data points. When exposed to new unseen Validation folds, this creates large extrapolation and numerical-instability errors on held-out folds, explaining the massive spikes observed in the right-tails of the CV curves. By strictly trusting the Out-Of-Fold Validation MSE to natively cap the models at Degree 3 and Degree 8, this methodology successfully avoided these catastrophic numerical illusions without relying on external Ridge or Lasso constraints.

\section{Conclusion}
The experimental results demonstrate that rigorous cross-validation and inclusive feature retention successfully produce robust unpenalized polynomial models. By focusing on explicit Out-of-Fold variance tracking rather than isolated Training metrics, Phase 1 and Phase 2 achieved out-of-fold $R^2$ estimates of $0.9067$ and $0.9947$ respectively, with the selected degrees determined by the evaluated cross-validation results.
\end{document}
